% ============================================================
%  ch04d：§4.7 结论（含定理 4.7.2.1 及证明）与参考文献
%  （原文 pp.39 中段–42；起于 §4.7，止于参考文献表与作者地址）
% ============================================================
\subsection{结论}

我们现在回到 §4.5 的设定。

\pnum{令 $\LS^{\arthm}_{\bT,\delta_{\bT}\neq0}\subseteq\LS^{\arthm}_{\bT}$ 表示 $\delta_{\bT}$ 的非零轨迹。}

我们有有理函数：
\[
  \LS^{\arthm,\circ}_{\bT}
  \stackrel{\text{定理 4.3.3.1}}{=}
  \bT^{\ad}/\bT=\bT\times^{\mathbf{B}}\bT
  \xrightarrow{p_1}\bT
  \xrightarrow{\zeta_{X,\bN,\bT}}\A^1,
  % [Tutor Note: 原文此处排作 \times^{\mathbf{B}}（粗体分类叠记号），疑为 \times_{\mathbf{B}\check T} 的排版变体；照录并登记。]
\]
它在 $\LS^{\arthm}_{\bT,\delta_{\bT}\neq0}$ 上显然有定义且可逆。按滥用记号，我们也用 $\zeta_{X,\bN,\bT}$ 记所得映射
\[
  \zeta_{X,\bN,\bT}\colon\LS^{\arthm}_{\bT,\delta_{\bT}\neq0}\to\A^1\setminus0.
\]

\pnum{我们现在证明以下结果：}

\begin{thm}
存在一个交换图：
{\scriptsize
\[
\begin{tikzcd}[column sep=3.5em, row sep=3.2em]
  \Gamma(\LS^{\arthm}_{\bT,\delta_{\bT}\neq0},\omega) \arrow[r, equal]
    & \Gamma(\LS^{\arthm}_{\bT},\omega)[\delta_{\bT}^{-1}] \arrow[rr, "\Eisspec"]
    & & \Gamma((\bG^{\ad}/\bG)^{\nonres},\omega)\\
  \Gamma(\LS^{\arthm}_{\bT,\delta_{\bT}\neq0},\omega) \arrow[r, equal]
    \arrow[u, "\cong"', "\zeta_{X,\bN,\bT}\cdot-"]
    & \Gamma(\LS^{\arthm}_{\bT},\omega)[\delta_{\bT}^{-1}] \arrow[r, "\mathrm{Eis}^{\mathrm{spec},\mathrm{toy}}"']
    & \Gamma(\bG^{\ad}/\bG,\omega)[(\delta_{\bG}\circ\tau_{\bG})^{-1}]
      \arrow[r, equal] \arrow[ur, "id"']
    & \Gamma((\bG^{\ad}/\bG)^{\nonres},\omega)
\end{tikzcd}
\]
}
\end{thm}

由推论 3.7.2.1，这显然给出 §4.5.6 中所需的满射性。因此，只剩下证明这个定理。

\begin{proof}[定理 4.7.2.1 的证明]
\textbf{第 1 步。}由紧合态射连接的叠的如下交换图
\[
\begin{tikzcd}[column sep=2.6em, row sep=1.8em]
  \mathbf{B}\bB \arrow[r, "\tau_{\bB}"] \arrow[d, "p^{\mathrm{toy}}"]
    & \LS^{\restr}_{\bB} \arrow[d, "p"]\\
  \mathbf{B}\bG \arrow[r, "\tau_{\bG}"]
    & \LS^{\restr}_{\bG}
\end{tikzcd}
\]
给出所得两个函子
\[
  \Rep(\bB)\to\IndCoh(\LS^{\restr}_{\bG})
\]
的一个等同，且保持紧对象。

对 Frobenius 取迹，这给出交换图
\begin{equation}
\begin{tikzcd}[column sep=3.2em, row sep=1.8em]
  \Gamma(\bB^{\ad}/\bB,\omega) \arrow[r] \arrow[d]
    & \Gamma(\LS^{\arthm}_{\bB},\omega) \arrow[d]\\
  \Gamma(\bG^{\ad}/\bG,\omega) \arrow[r]
    & \Gamma(\LS^{\arthm}_{\bG},\omega).
\end{tikzcd}
\end{equation}

\textbf{第 2 步。}其次，构造交换方块：
\[
\begin{tikzcd}[column sep=2.6em, row sep=1.8em]
  \mathbf{B}\bB \arrow[r, "\tau_{\bB}"] \arrow[d, "q^{\mathrm{toy}}"]
    & \LS^{\restr}_{\bB} \arrow[d, "q"]\\
  \mathbf{B}\bT \arrow[r, "\tau_{\bT}"]
    & \LS^{\restr}_{\bT}.
\end{tikzcd}
\]
我们得到自然变换
\[
  q^{*,\IndCoh}\tau_{\bT,*}^{\IndCoh}
  \to \tau_{\bB,*}^{\IndCoh}q^{\mathrm{toy},*,\IndCoh},
\]
它是函子
\[
  \Rep(\bT)\to\IndCoh(\LS^{\restr}_{\bB})
\]
之间的自然变换。但这个函子\emph{不是}同构，因此仅凭迹的函子性对它几乎无话可说。

尽管如此，我们断言有交换图：
\begin{equation}
\begin{tikzcd}[column sep=3.95em, row sep=4.4em, cells={font=\small}]
  \Gamma(\bT^{\ad}/\bT,\omega) \arrow[r, "p_{X,\bN,\bT}\cdot-"] \arrow[d, equal]
    & \Gamma(\bT^{\ad}/\bT,\omega) \arrow[r, "\tr(q^{*,\IndCoh}\tau_{\bT,*}^{\IndCoh})"]
    & \Gamma(\LS^{\arthm}_{\bB},\omega)\\
  \Gamma(\bT^{\ad}/\bT,\omega) \arrow[r, "q_{X,\bN,\bT}\cdot-"]
    & \Gamma(\bT^{\ad}/\bT,\omega) \arrow[ur, "\tr(\tau_{\bB,*}^{\IndCoh}q^{\mathrm{toy},*,\IndCoh})"']
    & 
\end{tikzcd}
\end{equation}

为构造这个图，注意这些映射自然是 $\Gamma(\LS^{\arthm}_{\bT},\OO)$-模的态射。这个代数显然通过其因子
\[
  \Gamma(\LS^{\arthm,\circ}_{\bT},\OO)=\Gamma(\bT^{\ad}/\bT,\OO)
\]
作用在 $\Gamma(\bT^{\ad}/\bT,\omega)$（我们想构造的图的源）上。因此，只需给出交换图
\[
\begin{tikzcd}[column sep=3.95em, row sep=4.4em, cells={font=\small}]
  \Gamma(\bT^{\ad}/\bT,\omega) \arrow[r, "p_{X,\bN,\bT}\cdot-"] \arrow[d, equal]
    & \Gamma(\bT^{\ad}/\bT,\omega) \arrow[r, "\tr(q^{*,\IndCoh}\tau_{\bT,*}^{\IndCoh})"]
    & \Gamma(\LS^{\arthm}_{\bB},\omega) \arrow[r, "\mathrm{proj}"] \arrow[d, equal]
    & \Gamma(\LS^{\arthm,\circ}_{\bB},\omega) \arrow[d, equal]\\
  \Gamma(\bT^{\ad}/\bT,\omega) \arrow[r, "q_{X,\bN,\bT}\cdot-"]
    & \Gamma(\bT^{\ad}/\bT,\omega) \arrow[r, "\tr(\tau_{\bB,*}^{\IndCoh}q^{\mathrm{toy},*,\IndCoh})"]
    & \Gamma(\LS^{\arthm}_{\bB},\omega) \arrow[r, "\mathrm{proj}"]
    & \Gamma(\LS^{\arthm,\circ}_{\bB},\omega),
\end{tikzcd}
\]
它是 $\Gamma(\bT^{\ad}/\bT,\OO)$-模的交换图（因为
$\LS^{\arthm,\circ}_{\bB}=\LS^{\arthm}_{\bB}\times_{\LS^{\arthm}_{\bT}}\LS^{\arthm,\circ}_{\bT}$）。

回忆 $H^{\ad}/H$ 上的对偶化层典范平凡。这个平凡化可以如下构造：等价 $\QCoh(\mathbf{B}H)\simeq\IndCoh(\mathbf{B}H)$ 在 Hochschild 同调上给出同构
$\Gamma(H^{\ad}/H,\OO)\simeq\Gamma(H^{\ad}/H,\omega)$。令
$\vol_H\in\Om^\infty\Gamma(H^{\ad}/H,\omega)$ 表示所得的生成元——显式地，它是 $\Rep(H)$ 的单位对象的类。因此，为得到上面的图，只需给出同构\footnote{技术上，为使 $p_{X,\bN,\bT}$ 与 $q_{X,\bN,\bT}$ 作用在这个等式的外侧（而不是作用在 $\vol_{\bT}$ 本身上），我们把所见的对象都看作上述意义上的 $\Gamma(\bT^{\ad}/\bT,\OO)$-模。}
\[
\begin{gathered}
  p_{X,\bN,\bT}\cdot\mathrm{proj}
  \Big(\tr\big(q^{*,\IndCoh}\tau_{\bT,*}^{\IndCoh}\big)(\vol_{\bT})\Big)\\
  =q_{X,\bN,\bT}\cdot\mathrm{proj}
  \Big(\tr\big(\tau_{\bB,*}^{\IndCoh}q^{\mathrm{toy},*,\IndCoh}\big)(\vol_{\bT})\Big)
  \in\Om^\infty\Gamma(\LS^{\arthm,\circ}_{\bB},\omega).
\end{gathered}
\]
由构造，我们有：
\[
  \tr\big(\tau_{\bB,*}^{\IndCoh}q^{\mathrm{toy},*,\IndCoh}\big)(\vol_{\bT})
  =\cl\big(\tau_{\bB,*}^{\IndCoh}\OO_{\mathbf{B}\bB},\alpha\big)
  \in\Om^\infty\Gamma(\LS^{\arthm}_{\bB},\omega).
\]
类似地，由基变换，我们有：
\[
  \tr\big(q^{*,\IndCoh}\tau_{\bT,*}^{\IndCoh}\big)(\vol_{\bT})
  =\cl\big(\OO_{\LS^{\restr}_{\bB}\times_{\LS^{\restr}_{\bT}}\mathbf{B}\bT},\alpha\big)
  =\cl\big(\OO_{\LS^{\restr}_{\bB;\bT}},\alpha\big)
  \in\Om^\infty\Gamma(\LS^{\arthm}_{\bB},\omega).
\]
于是该恒等式由引理 4.6.6.4 推出，注意
$\LS^{\arthm}_{\bB;\bT}=\LS^{\arthm,\circ}_{\bB}$（见注 4.6.6.3）。

\textbf{第 3 步。}拼接图 (4.16) 与 (4.17)，我们得到交换图：
\[
\begin{tikzcd}[column sep=1.1em, row sep=3.4em]
  \Gamma(\bT^{\ad}/\bT,\omega) \arrow[r, "p_{X,\bN,\bT}\cdot-"] \arrow[d, equal]
    & \Gamma(\bT^{\ad}/\bT,\omega) \arrow[r, "\tr(\tau_{\bT,*}^{\IndCoh})"]
    & \Gamma(\LS^{\arthm}_{\bT},\omega) \arrow[r, "\Eisspec"]
    & \Gamma(\LS^{\arthm}_{\bG},\omega) \arrow[d, equal]\\
  \Gamma(\bT^{\ad}/\bT,\omega) \arrow[r, "q_{X,\bN,\bT}\cdot-"]
    & \Gamma(\bT^{\ad}/\bT,\omega) \arrow[r, "\mathrm{Eis}^{\mathrm{spec},\mathrm{toy}}"']
    & \Gamma(\bG^{\ad}/\bG,\omega) \arrow[r, "\tr(\tau_{\bG,*}^{\IndCoh})"]
    & \Gamma(\LS^{\arthm}_{\bG},\omega).
\end{tikzcd}
\]
这个图细化了我们本该证明的定理。
\end{proof}

\begin{thebibliography}{99}
\bibitem{arinkin-gaitsgory} D. Arinkin and D. Gaitsgory, \emph{Singular support of coherent sheaves and the geometric Langlands conjecture}, Selecta Math. \textbf{21} (2015), no. 1, 1--199. Springer.
\bibitem{AGKRV1} D. Arinkin, D. Gaitsgory, D. Kazhdan, S. Raskin, N. Rozenblyum, and Y. Varshavsky, \emph{The stack of local systems with restricted variation and geometric Langlands theory with nilpotent singular support}, preprint, 2020. arXiv:2010.01906.
\bibitem{AGKRV2} D. Arinkin, D. Gaitsgory, D. Kazhdan, S. Raskin, N. Rozenblyum, and Y. Varshavsky, \emph{Duality for automorphic sheaves with nilpotent singular support}, preprint, 2020. arXiv:2012.07665.
\bibitem{AGKRV3} D. Arinkin, D. Gaitsgory, D. Kazhdan, S. Raskin, N. Rozenblyum, and Y. Varshavsky, \emph{Automorphic functions as the trace of Frobenius}, preprint, 2021. arXiv:2102.07906.
\bibitem{beilinson} A. Beilinson, \emph{Constructible sheaves are holonomic}, Selecta Math. \textbf{22} (2016), no. 4, 1797--1819. Springer.
\bibitem{braverman-gaitsgory} A. Braverman and D. Gaitsgory, \emph{Deformations of Local Systems and Eisenstein Series}, Geom. Funct. Anal., 2008. Springer Science+Business Media.
\bibitem{benzvi-nadler} D. Ben-Zvi and D. Nadler, \emph{Betti geometric Langlands}, in Proceedings of Symposia in Pure Mathematics, vol. 97, 2018, pp. 3--41.
\bibitem{drinfeld-gaitsgory1} V. Drinfeld and D. Gaitsgory, \emph{On some finiteness questions for algebraic stacks}, Geom. Funct. Anal. \textbf{23} (2013), no. 1, 149--294. Springer.
\bibitem{drinfeld-gaitsgory2} V. Drinfeld and D. Gaitsgory, \emph{Geometric constant term functor(s)}, Selecta Math. \textbf{22} (2016), no. 4, 1881--1951. Springer.
\bibitem{faergeman-raskin} J. Faergeman and S. Raskin, \emph{Non-vanishing of geometric Whittaker coefficients for reductive groups}, J. Amer. Math. Soc., 2025.
\bibitem{fargues-scholze} L. Fargues and P. Scholze, \emph{Geometrization of the local Langlands correspondence}, preprint, 2021. arXiv:2102.13459.
\bibitem{gaitsgory-indcoh} D. Gaitsgory, \emph{Ind-coherent sheaves}, Mosc. Math. J. \textbf{13} (2013), no. 3, 399--528.
\bibitem{gaitsgory-outline} D. Gaitsgory, \emph{Outline of the proof of the geometric Langlands conjecture for $\mathrm{GL}(2)$}, Ast\'erisque, 2015. Centre National de la Recherche Scientifique.
\bibitem{gaitsgory-sheaves} D. Gaitsgory, \emph{Sheaves of categories and the notion of 1-affineness}, in Stacks and categories in geometry, topology, and algebra, vol. 643, 2015, pp. 127--225. American Mathematical Society, Providence, RI.
\bibitem{GKRV} D. Gaitsgory, D. Kazhdan, N. Rozenblyum, and Y. Varshavsky, \emph{A toy model for the Drinfeld--Lafforgue shtuka construction}, Indag. Math., 2021. Elsevier.
\bibitem{gaitsgory-raskin} D. Gaitsgory and S. Raskin, \emph{Proof of the geometric Langlands conjecture I: construction of the functor}, arXiv preprint arXiv:2405.03599, 2024.
\bibitem{gaitsgory-rozenblyum1} D. Gaitsgory and N. Rozenblyum, \emph{A study in derived algebraic geometry: Volume I: correspondences and duality}, Mathematical Surveys and Monographs, vol. 221, American Mathematical Society, 2019.
\bibitem{gaitsgory-rozenblyum2} D. Gaitsgory and N. Rozenblyum, \emph{A study in derived algebraic geometry: Volume II: deformations, Lie theory and formal geometry}, Mathematical Surveys and Monographs, vol. 221, American Mathematical Society, 2020.
\bibitem{harder} G. Harder, \emph{\"Uber die Galoiskohomologie halbeinfacher algebraischer Gruppen. III}, Walter de Gruyter, Berlin/New York, 1975.
\bibitem{lafforgue-l} L. Lafforgue, \emph{Chtoucas de Drinfeld et correspondance de Langlands}, Invent. Math. \textbf{147} (2002), no. 1, 1--241. Springer.
\bibitem{lafforgue-v1} V. Lafforgue, \emph{Quelques calculs reli\'es \`a la correspondance de Langlands g\'eom\'etrique pour $\mathbf{P}^1$}, preprint, 2009. Available at \url{http://www.math.jussieu.fr/~vlafforg}.
\bibitem{lafforgue-v2} V. Lafforgue, \emph{Chtoucas pour les groupes r\'eductifs et param\'etrisation de Langlands globale}, J. Amer. Math. Soc. \textbf{31} (2018), no. 3, 719--891.
\bibitem{laumon1} G. Laumon, \emph{Correspondance de Langlands g\'eom\'etrique pour les corps de fonctions}, Duke Math. J. \textbf{54} (1987), no. 2, 309--359. Duke University Press.
\bibitem{laumon2} G. Laumon, \emph{Transformation de Fourier g\'en\'eralis\'ee}, arXiv preprint alg-geom/9603004, 1996.
\bibitem{lurie} J. Lurie, \emph{Spectral algebraic geometry}, preprint, 2018. Available at \url{https://www.math.ias.edu/~lurie/papers/moduli.pdf}.
\bibitem{moeglin-waldspurger} C. Moeglin and J.-L. Waldspurger, \emph{Spectral decomposition and Eisenstein series}, Cambridge University Press, 1995. London Mathematical Society Monographs, Number 113.
\bibitem{nadler-yun} D. Nadler and Z. Yun, \emph{Spectral action in Betti Geometric Langlands}, Israel J. Math. \textbf{232} (2019), no. 1.
\bibitem{deligne} P. Deligne, \emph{S\'eminaire de G\'eom\'etrie Alg\'ebrique du Bois Marie--Cohomologie \'etale--(SGA 4 1/2)}, Springer, Berlin/New York, 1977.
\bibitem{varshavsky} Y. Varshavsky, \emph{Moduli spaces of principal F-bundles}, Selecta Math. (N.S.) \textbf{10} (2004), no. 1, 131. Springer.
\bibitem{xue-finiteness} C. Xue, \emph{Finiteness of cohomology groups of stacks of shtukas as modules over Hecke algebras, and applications}, \'Epijournal de G\'eom\'etrie Alg\'ebrique \textbf{4}, 2020. Episciences.org.
\bibitem{xue-smoothness} C. Xue, \emph{Smoothness of cohomology sheaves of stacks of shtukas}, preprint, 2020. arXiv:2012.12833.
\bibitem{zhu} X. Zhu, \emph{Coherent sheaves on the stack of Langlands parameters}, preprint, 2020. arXiv:2008.02998.

\medskip
\noindent Yale University, Department of Mathematics, 219 Prospect St, New Haven, CT 06511\\
\noindent\emph{Email address:} \texttt{sam.raskin@yale.edu}
\end{thebibliography}
